% Figure: the nth roots of unity for n = 6, sitting at the corners of a
% regular hexagon inscribed in the unit circle, with the cube roots of
% unity highlighted as a subset.
\begin{figure}[htbp]
\centering
\begin{tikzpicture}[scale=2.2, line cap=round, >=latex]
  \draw[->, engblack, thin] (-1.5, 0) -- (1.6, 0) node[right, font=\small] {$\Re$};
  \draw[->, engblack, thin] (0, -1.5) -- (0, 1.6) node[above, font=\small] {$\Im$};
  \draw[engblack!40, thin, dashed] (0, 0) circle (1);
  % Six sixth-roots of unity at multiples of 60 degrees.
  \foreach \k in {0,1,2,3,4,5} {
    \pgfmathsetmacro{\ang}{\k*60}
    \filldraw[engblue] (\ang:1) circle (0.035);
    \draw[engblue!50, thin] (0,0) -- (\ang:1);
  }
  % Hexagon edges.
  \draw[engblue, thick]
    (0:1) -- (60:1) -- (120:1) -- (180:1) -- (240:1) -- (300:1) -- cycle;
  % Highlight the three cube roots of unity (multiples of 120 deg).
  \foreach \k in {0,1,2} {
    \pgfmathsetmacro{\ang}{\k*120}
    \draw[engred, thick, ->] (0,0) -- (\ang:1);
    \filldraw[engred] (\ang:1) circle (0.045);
  }
  \node[engred, font=\footnotesize, anchor=west] at (0:1.05) {$1$};
  \node[engred, font=\footnotesize, anchor=south] at (120:1.08) {$e^{i 2\pi/3}$};
  \node[engred, font=\footnotesize, anchor=north] at (240:1.08) {$e^{i 4\pi/3}$};
  \node[engblue, font=\footnotesize, anchor=south] at (60:1.10) {$e^{i\pi/3}$};
  \node[engblue, font=\footnotesize, anchor=north] at (300:1.10) {$e^{i 5\pi/3}$};
\end{tikzpicture}
\caption[Roots of unity as a regular polygon]{The six sixth-roots of
unity (blue) sit at the corners of a regular hexagon inscribed in the
unit circle, spaced by $2\pi/6 = 60^{\circ}$. The three cube roots of
unity (red) are the subset at $0$, $2\pi/3$, and $4\pi/3$; they are the
phasors of a balanced three-phase system and they sum to zero. The
$n$th roots of unity are the working basis of the length-$n$ discrete
Fourier transform and of the $n$-phase electrical system. Each root is
$e^{i 2\pi k/n}$ for $k = 0, \ldots, n-1$, and the sum of all $n$ roots
of unity is zero whenever $n \ge 2$.}
\label{fig:vol02:ch03:roots-unity}
\end{figure}
